
\begin{proposition}[K3 components in arbitrary torus fillings are exactly the cyclic generalized Kummer surfaces]\label{cdp:st:C.2.UN.32}
The K3 minimal birational classes occurring as
normalization components of proper two-torus
degenerations are exactly the minimal resolutions
of cyclic quotients of complex two-tori by
nontrivial symplectic automorphisms.
The generator has no eigenvalue one on the tangent space
and, after an actual translation of the origin,
is linear of order $2,3,4$, or $6$.
The exceptional configurations are those of
(C.K3.1). This criterion holds in every
algebraic dimension.

Every such class occurs in a smooth compact
filling with both special reduced supports
of Euler number zero, by the explicit quotient
and balancing construction of \ref{cdp:st:C.2.UN.24}.
In particular the algebraic-dimension-zero
equivalence of that statement extends to
all K3 components.
\end{proposition}
\begin{proof}
Let $M$ be the K3 minimal model of an original
component. Take the actual multiplicity-one
component cover $E$ and its cyclic stabilizer
$G$ from \ref{cdp:st:C.2.UN.16}. The proper graph comparisons
of \ref{cdp:st:C.2.UN.19} give
\[
                         E/G\dashrightarrow M
                                                               \tag{C.KX.1}\label{cdp:eq:C.KX.1}
\]
as a proper birational comparison after
resolving the graph. Pullback of the K3
two-form is injective, whereas the
compact-core bound \ref{cdp:st:C.2.UN.28} gives $p_g(E)\le1$.
Thus
\[
                       p_g(E)=1,\qquad\kappa(E)\ge0.
                                                               \tag{C.KX.2}\label{cdp:eq:C.KX.2}
\]
Pass to its unique smooth minimal model
$E_0$; the finite action descends through
the actual point blowdowns.
If $E_0$ is a torus, \ref{cdp:st:C.2.UN.23} already gives
the required cyclic K3 quotient.

In the other cases, \ref{cdp:st:C.2.UN.22} and \ref{cdp:st:C.2.UN.29} give
an actual relatively minimal elliptic map
\[
                   q:E_0\longrightarrow B
                                                               \tag{C.KX.3}\label{cdp:eq:C.KX.3}
\]
with identity integral monodromy and
only smooth multiple fibres. Its
relative Jacobian is one fixed elliptic
curve $F$. The local invariant
direct-image calculation (C.RG.13)
applies and gives
$p_g(E_0)=g(B)$, so $B$ is elliptic.
The map can be taken $G$-equivariant.
In the rank-one branch this is \ref{cdp:st:C.2.UN.20}.
In algebraic dimension one its
connected-fibre function field is
intrinsic, by (C.RG.13) and
(C.K3.4). In algebraic dimension
two a non-torus minimal model
with $p_g=1$ and $\chi(\OO)=0$
has Kodaira dimension one, so
its elliptic map is its intrinsic
Iitaka fibration, as in \ref{cdp:st:C.2.UN.24}.
Algebraic dimension zero was
already the torus branch of \ref{cdp:st:C.2.UN.16}.
These are the same equivariant
graph and Stein comparisons
used there.

We first describe every multiple
fibre with its actual cyclic
uniformization. For a fibre $mC$,
the normalized root change
$t=u^m$ is \'etale over the
original smooth surface and
has connected smooth elliptic
central fibre. The identity
monodromy and constant period
ratio make this root family
a torsor under the fixed $F$
over a disc. Its torsor is
trivial: the sheaf sequence
$0\to\Lambda_F\to\OO\to F\to0$
on a disc has $H^1(\OO)=0$
and $H^2(\Lambda_F)=0$.
Thus it is an actual product.
The generator of its root
deck group has translation
linear part and can be
written
$(u,z)\mapsto(\zeta_m u,z+\tau+a(u))$,
where $\tau=a_{\rm total}(0)$
and $a(0)=0$.
Lift $a$ holomorphically to
$\C$. The finite-order
identity gives
$\sum_{j=0}^{m-1}a(\zeta_m^ju)=0$.
Replacing $z$ by $z+b(u)$,
with
$b(u)=-m^{-1}\sum_{j=0}^{m-1}j\,a(\zeta_m^ju)$,
removes $a(u)$.
The local model is therefore
\[
 (F\times\Delta_u)/\langle\delta\rangle,
 \qquad
 \delta(u,z)=(\zeta_m u,z+\tau),\quad
 \operatorname{ord}(\tau)=m,\quad t=u^m .
                                                               \tag{C.KX.4}\label{cdp:eq:C.KX.4}
\]
The exact order of $\tau$
follows from freeness of
the \'etale root covering
on its central fibre.

Let $H$ be the kernel of the
action of $G$ on $B$.
Every element of $H$ acts
on a general elliptic fibre
by a translation. Indeed
a nonidentity linear part
would make the general
finite quotient fibre
rational. The actual
quotient and Stein map
would then give a rational
fibration on a birational
model of $M$, contradicting
its nonzero canonical form
by restriction to a general
rational fibre.
The translation subgroup
on the fixed $F$ is
constant: its points have
bounded finite order and
belong to a discrete
finite torsion set.

The quotient
\[
              E_1=E_0/H,\qquad Q=G/H
                                                               \tag{C.KX.5}\label{cdp:eq:C.KX.5}
\]
is smooth. To check this
also at multiple fibres,
use \eqref{cdp:eq:C.KX.4}. A fibrewise
finite automorphism has
linear part one on the
reduced elliptic fibre,
so its restriction is
a translation. If it
fixes one point, it
fixes that curve pointwise.
Finite holomorphic
linearization then gives
$(x,z)\mapsto(\zeta_e x,z)$;
its quotient coordinates
are $(x^e,z)$. Every
point stabilizer is
cyclic, so these charts
prove smoothness of the
full quotient. They also
show exactly how the
fibre multiplicity can
decrease: $t=x^m$ becomes
$t=(x^e)^{m/e}$, with
$e\mid m$. The reduced
support remains a smooth
elliptic curve. The
general fibre is the
fixed isogenous elliptic
curve obtained by the
constant translation
subgroup. Its integral
linear monodromy remains
the identity.

Thus $q_1:E_1\to B$ is
again a smooth-multiple
elliptic fibration with
constant Jacobian.
The same direct-image
calculation gives
\[
 p_g(E_1)=1,\quad q(E_1)=2,\quad
                 \chi(\OO_{E_1})=0 .
                                                               \tag{C.KX.6}\label{cdp:eq:C.KX.6}
\]
The actual quotient
$E_1/Q=E_0/G$ remains
birational to $M$.
The group $Q$ acts
faithfully on the base.
Its quotient fibration
has connected elliptic
general fibre. The
K3 descent argument
(C.K3.5) makes its
base $B/Q\simeq\PP^1$.
Consequently $Q$ is
nontrivial and its
action on $B$ is not
a group of translations.
After translating the
origin on $B$, its
generator is a linear
elliptic automorphism.
The derivative character
$\lambda:Q\to\C^*$ is
faithful.

We record the exact
canonical character.
Choose a nonzero
one-form $\omega_B$
on the elliptic base
and a time form $dz$
for the fixed Jacobian.
Their wedge glues
over the translation
trivializations. In
\eqref{cdp:eq:C.KX.4} it is
\[
 \Omega=\omega_B\wedge dz
        =m\,c(u^m)u^{m-1}\,du\wedge dz ,
                  \qquad c(0)\ne0 .
                                                               \tag{C.KX.7}\label{cdp:eq:C.KX.7}
\]
It is invariant under
$\delta$ and extends
holomorphically. Its
divisor is exactly
$\sum_b(m_b-1)C_b$.
The group acts on
the fixed elliptic
Jacobian with a
constant linear
character $\beta$;
translation terms
add only base
differentials to $dz$.
Hence
$g^*\Omega=\lambda(g)\beta(g)\Omega$.
The nonzero K3 form
pulls back through
the actual graph to
the unique canonical
form on $E_1$.
It is invariant, so
\[
                   \lambda(g)\beta(g)=1
                       \quad(g\in Q).                \tag{C.KX.8}\label{cdp:eq:C.KX.8}
\]

Here is the local
multiple-fibre
obstruction, with
all covering maps
retained. Let a
base point have
stabilizer generated
by $g$ of order
$s>1$, and suppose
its fibre has
multiplicity $m$.
In \eqref{cdp:eq:C.KX.4} choose
a lift of $g$:
\[
 \widetilde g(u,z)
       =(\rho u,\beta z+a(u)),\qquad
       \rho^m=\lambda .
                                                               \tag{C.KX.9}\label{cdp:eq:C.KX.9}
\]
Such a lift is the
functorial lift of
the automorphism
to the normalized
root change, after
choosing an $m$th
root of $\lambda$.
Conjugation preserves
the root deck group.
Its action on $u$
shows that
$\widetilde g\delta
 \widetilde g^{-1}=\delta$.
Restricting this
identity to $u=0$
gives
\[
                         \beta\tau=\tau .
                                                               \tag{C.KX.10}\label{cdp:eq:C.KX.10}
\]
By \eqref{cdp:eq:C.KX.8}, $\beta$
is nonidentity and
has order $s$.
The sum of its
endomorphisms
$1+\beta+\cdots+\beta^{s-1}$
is zero. Applying
it to $\tau$ in
\eqref{cdp:eq:C.KX.10} gives
$s\tau=0$.
Since $\tau$ has
exact order $m$,
we obtain $m\mid s$.

On the reduced
elliptic fibre
$F/\langle\tau\rangle$,
the automorphism
$g$ has nonidentity
linear part $\beta$.
It has a fixed
point, because
$1-\beta$ is an
isogeny of that
elliptic curve.
Choose a lift of
this point and
multiply the
chosen $\widetilde g$
by a root deck
transformation so
that it fixes
that lift.
Then $\widetilde g^s$
is a root deck
transformation
fixing that point.
The deck action
is free, so
$\widetilde g^s=1$.
In particular
$\rho^s=1$, while
$\rho^m=\lambda$
has exact order
$s$. This implies
$\gcd(m,s)=1$.
Together with
$m\mid s$, it
proves
\[
          \text{a branched base value has }m=1 .
                                                               \tag{C.KX.11}\label{cdp:eq:C.KX.11}
\]

At a base point
with trivial
stabilizer, the
quotient is locally
an isomorphism
on a neighbourhood
of one fibre in
the orbit. Its
smooth elliptic
support therefore
retains multiplicity
$m$ and the zero
order $m-1$ of
\eqref{cdp:eq:C.KX.7}.
On a common smooth
resolution of
$E_1/Q$ and $M$,
the strict transform
of this curve is
elliptic and cannot
be contracted by
a birational map
of smooth surfaces:
all curves
exceptional for
successive point
blowdowns are
rational.
At its generic
point the pullback
of the nowhere-zero
K3 form must
therefore have
zero order zero.
Its equality with
the descended form
\eqref{cdp:eq:C.KX.7} forces
$m-1=0$.
This treats every
free base orbit.
We have proved
that $q_1$ has
no multiple fibre
at all.

We finish by an
explicit Chern-class
comparison, without
assuming that a
finite image or
preimage is
K\"ahler.
The smooth map
$q_1:E_1\to B$
is now a principal
bundle under its
fixed elliptic
curve $F_1$.
Its Chern vector is
\[
        c\in H^2(B,\Lambda_{F_1})
                  \simeq\Lambda_{F_1}.
                                                               \tag{C.KX.12}\label{cdp:eq:C.KX.12}
\]
A bundle
automorphism
covering $g_B$
and having
fibre linear
part $\beta$
gives
$g_B^*c=\beta_*c$.
For the precise
map, lift the
translation
transition
functions on
a base cover
to $\C$.
Their triple
differences are
the lattice
cocycle for $c$.
Conjugating the
trivializations
multiplies that
cocycle by $\beta$,
up to a lattice
coboundary; this
is the displayed
equality.
The holomorphic
automorphism of
the elliptic
base has degree
one, so it acts
as the identity
on $H^2(B,\Z)$.
Thus $c=\beta c$.
For a generator,
$\beta\ne1$, and
$\beta-I$ is
invertible on
$\Lambda_{F_1}
\otimes\R$.
The lattice is
torsion-free, so
$c=0$.

The sheaf sequence
\[
 0\longrightarrow\Lambda_{F_1}
 \longrightarrow\OO_B
 \longrightarrow\mathcal F_1
 \longrightarrow0
                                                               \tag{C.KX.13}\label{cdp:eq:C.KX.13}
\]
now lifts the
principal-bundle
class to
$H^1(B,\OO_B)$.
Write $B=\C/\Gamma$.
Every class in
that one-dimensional
complex space
is represented
by a constant
$(0,1)$-form
$a\,d\overline z$.
On the universal
cover its
primitive is
$a\overline z$,
and its changes
under $\gamma\in\Gamma$
are the constants
$a\overline\gamma$.
Consequently its
additive torsor,
and its image
under \eqref{cdp:eq:C.KX.13},
have the exact
quotient presentation
\[
 E_1\simeq
 \C^2\big/
 \left\langle
 (z,w)\mapsto
       (z+\gamma,w+a\overline\gamma),\
 (z,w)\mapsto(z,w+\lambda)
 \right\rangle_{
       \gamma\in\Gamma,\,
       \lambda\in\Lambda_{F_1}} .
                                                               \tag{C.KX.14}\label{cdp:eq:C.KX.14}
\]
Equality of the
Cech classes gives
this actual
isomorphism:
the difference
of two transition
systems is a
holomorphic
zero-cochain,
and translating
the fibre
coordinate by
that cochain
glues the local
maps.
The displayed
group is a
discrete rank-four
lattice of
translations,
with compact
fundamental
parallelepiped.
Thus $E_1$
is a complex
two-torus.

Finally $Q$ is
cyclic and
$E_1/Q=E_0/G$
has the actual
birational
comparison
\eqref{cdp:eq:C.KX.1} to $M$.
\ref{cdp:st:C.2.UN.23} applies
to this torus,
and gives
exactly its
symplectic K3
quotient case.
The order and
exceptional
configuration
calculation
(C.K3.6)--(C.K3.7)
gives $2,3,4,6$
and (C.K3.1).
Conversely every
one of those
actual torus
quotients is
realized by
(C.K3.10), its
proper resolution,
and the two
explicit balancing
operations of
\ref{cdp:st:C.2.UN.24}. This
proves both
directions in
every algebraic
dimension.
\end{proof}


